\label{anexo:repro}

Este anexo reúne la información necesaria para reproducir íntegramente el experimento descrito en el capítulo de evaluación. Se incluye porque una evaluación cuyos resultados no pueden rehacerse por un tercero aporta poca evidencia, con independencia de lo cuidadoso que haya sido su diseño.

\section{Entorno y dependencias}
El prototipo se implementó con nivel de lenguaje Java~17 y se verificó sobre JDK~17 y JDK~26. La única dependencia externa es \textit{ArchUnit} en su versión 1.3.0, empleada exclusivamente para importar el \textit{bytecode} y resolver las dependencias entre clases; el resto de la lógica —el algoritmo de detección de ciclos, el motor de conformidad y los exportadores— no depende de bibliotecas de terceros.

La biblioteca se distribuye junto al repositorio, de modo que la compilación y la ejecución no requieren Maven, repositorio de artefactos ni conexión a internet. Esta decisión responde al requisito de ligereza (RNF2) y, de paso, hace que el experimento siga siendo reproducible aunque los repositorios remotos cambien. Todas las operaciones se invocan mediante un único lanzador, \texttt{arquitecta.bat} en Windows o \texttt{./arquitecta.sh} en macOS y Linux, que compila el prototipo la primera vez y cada vez que su código cambia.

\section{Estructura del repositorio}
\begin{verbatim}
prototipo/
  arquitecta.bat / .sh        lanzador unico (aplicacion y tareas)
  src/main/java/              nucleo, servicio, interfaz y asistente
  src/test/java/              bateria de pruebas (272 comprobaciones)
  lib/                        ArchUnit 1.3.0 y slf4j-api
  caso-estudio/
    c0-base/                  escenario C0 (linea base)
    c1-capas/                 escenario C1 (ruptura de capas)
    c2-ciclo/                 escenario C2 (ciclo de dependencia)
    c3-concentracion/         escenario C3 (componente dios)
    c4-combinado/             escenario C4 (las tres derivas)
    rules.properties          catalogo de reglas del caso de estudio
  ejemplos/
    autoanalisis/             catalogo para el autoanalisis
    archunit/                 configuraciones de la seccion 4.6
    sistemas/                 catalogos y resultados de la seccion 4.7
    estudios/                 evolucion, umbral y oraculo (4.8 a 4.10)
    asistente/                exploracion del asistente (anexo B)
    graficos/                 grafos en SVG generados por la herramienta
    interfaz/                 capturas de la interfaz grafica
  scripts/                    guiones que invoca el lanzador
decisiones/                   registro de decisiones (ADR, anexo C)
\end{verbatim}

\noindent Cada escenario tiene la estructura de un proyecto Maven: su código está en \texttt{src/main/java} y la compilación deja sus clases en \texttt{target/classes}, de modo que la carpeta del escenario puede indicarse tal cual, desde la línea de comandos o desde la interfaz gráfica.

\section{Pasos para reproducir el experimento}
\begin{enumerate}
    \item \textbf{Compilar.} Ejecutar \texttt{arquitecta compilar} (\texttt{./arquitecta.sh compilar} en macOS o Linux). El guion compila el núcleo y la batería de pruebas en \texttt{out/}, empaqueta la herramienta en \texttt{out/arquitecta.jar} y compila cada escenario en su \texttt{target/classes}. El paso es opcional: los siguientes compilan por su cuenta si hace falta.
    \item \textbf{Verificar el núcleo.} Ejecutar \texttt{arquitecta pruebas}. La batería debe reportar 272 comprobaciones superadas y terminar con código de salida cero; cualquier fallo detiene la reproducción y señala el caso concreto que discrepa.
    \item \textbf{Ejecutar el experimento.} Ejecutar \texttt{arquitecta experimento}. El guion imprime el catálogo de reglas, recorre los cinco escenarios bajo los dos modos de clasificación y deposita la evidencia en \texttt{resultados/}.
    \item \textbf{Validación sobre sistemas reales (opcional).} Ejecutar primero \texttt{arquitecta aplicaciones}, que clona las tres aplicaciones en capas en el \textit{commit} fijado y las compila ---con Docker y la imagen oficial de Maven con JDK~17 si está disponible, o con el \textit{Maven Wrapper} de cada proyecto---, y después \texttt{arquitecta sistemas}. Este segundo guion descarga de Maven Central los \texttt{.jar} de las bibliotecas, verifica su suma SHA-1, inyecta las derivas descritas en la sección~\ref{sec:sistemas} y escribe la tabla de resultados y la ubicación de cada deriva en \texttt{ejemplos/sistemas/}. Si alguna aplicación no está compilada, se omite con un aviso. Requiere conexión la primera vez.
    \item \textbf{Estudios complementarios (opcional).} Ejecutar \texttt{arquitecta estudios}, que reproduce la evolución histórica (sección~\ref{sec:evolucion}), la calibración del umbral (sección~\ref{sec:umbral}), el contraste con JDepend (sección~\ref{sec:oraculo}) y la comparación con \textit{ArchUnit} (sección~\ref{sec:aporte}); cada uno puede ejecutarse por separado con \texttt{arquitecta estudios evolucion}, \texttt{umbral}, \texttt{oraculo} o \texttt{archunit}. Descarga y verifica las treinta y cinco versiones históricas y JDepend la primera vez, y escribe la evidencia en \texttt{ejemplos/estudios/}.
    \item \textbf{Contrastar.} Comparar el archivo \texttt{resultados/modo-cr/resumen.csv} con la tabla~\ref{tab:resultados} y el archivo \texttt{resultados/modo-cr-severidad/resumen.csv} con la tabla~\ref{tab:refinamiento}.
\end{enumerate}

Para reproducir el autoanálisis descrito en el capítulo de desarrollo basta con dirigir el mecanismo hacia su propio directorio de clases compiladas:

\begin{verbatim}
arquitecta analizar out\arquitecta-core ^
     ejemplos\autoanalisis\rules-arquitecta.properties
\end{verbatim}

\section{Figuras de grafos y sistema externo}
Las figuras de grafos de este documento se regeneran con \texttt{arquitecta figuras}, que escribe las versiones TikZ en el directorio \texttt{figuras/} de la monografía y las SVG en \texttt{ejemplos/graficos/} del prototipo. Para obtener solo el dibujo de un sistema cualquiera basta con indicar las opciones \texttt{-{}-svg} o \texttt{-{}-tikz}:

\begin{verbatim}
arquitecta analizar lib\archunit-1.3.0.jar ^
     ejemplos\archunit\rules-archunit.properties --svg archunit.svg
\end{verbatim}

\noindent El mismo análisis puede hacerse desde la interfaz gráfica, que se abre ejecutando el lanzador sin argumentos ---o con doble clic sobre él--- y permite elegir el \texttt{.jar} y el catálogo con un selector de archivos; su pantalla de inicio ofrece, además, los cinco escenarios del caso de estudio para analizarlos con un clic. El ejemplo reproduce además el análisis de la sección~\ref{sec:archunit}: el mecanismo admite directamente el archivo \texttt{.jar}, y el directorio \texttt{ejemplos/archunit/} contiene las configuraciones con y sin exenciones cuyos resultados recoge la tabla~\ref{tab:archunit}.

\section{Evidencia generada}
Cada ejecución con la opción \texttt{-{}-salida} produce cuatro artefactos por escenario y una fila en el resumen comparativo:

\begin{itemize}
    \item un \textbf{reporte legible} (\texttt{.txt}) con la Tasa de Conformidad, los estados deducidos por cada criterio, los hallazgos con su evidencia y la declaración de alcance del diagnóstico;
    \item la \textbf{evidencia estructurada} (\texttt{.json}), que incluye el grafo completo, los conteos por regla y la totalidad de las trazas de cada hallazgo, sin el recorte que el reporte legible aplica por brevedad;
    \item el \textbf{grafo de dependencias} dibujado en tres formatos con la misma disposición: una imagen \texttt{.svg} que se abre en cualquier navegador, una figura \texttt{.tikz} lista para incluirse en \LaTeX{} y un archivo \texttt{.dot} para quien prefiera procesarlo con Graphviz;
    \item los \textbf{hallazgos tabulados} (\texttt{.csv}), una fila por hallazgo, para su tratamiento en una hoja de cálculo.
\end{itemize}

El archivo \texttt{resumen.csv} de cada modo agrega una fila por escenario con el número de paquetes, dependencias, verificaciones e incumplimientos, la Tasa de Conformidad, los tres estados calculados y la CR por regla. Es el artefacto del que proceden directamente las tablas del capítulo de evaluación.

\section{Determinismo de los resultados}
El mecanismo no emplea aleatoriedad, concurrencia ni marcas temporales en su cálculo, por lo que ante la misma entrada produce siempre el mismo diagnóstico. La construcción del grafo recorre las clases y sus dependencias en orden alfabético, y las estructuras que lo recorren conservan el orden de inserción, de modo que también el orden de los hallazgos y de su evidencia es estable entre ejecuciones y no depende de cómo se nombre la ruta analizada ni del orden en que la biblioteca de importación entregue las clases. Ello permite comparar dos corridas mediante una simple diferencia textual de sus archivos de evidencia, sin necesidad de normalización previa. Las capturas de la interfaz gráfica son la única excepción: dependen de las fuentes tipográficas del equipo y muestran el tiempo de análisis, por lo que se regeneran junto con las figuras pero no se comparan byte a byte.

Al ejecutar sobre JDK~26 aparecen dos avisos en la salida de error —la ausencia de un proveedor de \textit{logging} para SLF4J y una advertencia sobre el uso de \texttt{sun.misc.Unsafe} emitida por una biblioteca embebida en ArchUnit— que no afectan al resultado y no se presentan con JDK~17.
